\documentclass[10pt,a4paper]{amsart}
\usepackage[margin=19mm]{geometry}
\usepackage{amsmath,amssymb,mathtools}
\usepackage[scr=rsfs,cal=euler]{mathalfa}
\usepackage{xfrac}
\usepackage[unicode,hidelinks,bookmarks=false]{hyperref}
\newcommand{\ie}{i.e.\ }
\newcommand{\cf}{cf.\ }
\newcommand{\resp}{respectively\ }
\newcommand{\Subsection}{\subsection}
\providecommand{\nobreakdash}{\nobreak\hbox{-}}
\setlength{\emergencystretch}{2em}
\multlinegap0pt
\usepackage{bm}
\usepackage{bbm}
\usepackage{dsfont}
\usepackage{xspace}
\usepackage{cancel}
\usepackage{mathtools}
\usepackage{amsthm}
\makeatletter
\@ifundefined{thm}{\newtheorem{thm}{Theorem}}{}
\@ifundefined{prop}{\newtheorem{prop}[thm]{Proposition}}{}
\makeatother
\newcommand{\NN}{{\mathbb N}}
\newcommand{\aA}{{\mathbf a}}
\newcommand{\bb}{{\mathbf b}}
\newcommand{\fh}{{\mathfrak h}}
\newcommand{\pp}{{\mathbf p}}
\newcommand{\qq}{{\mathbf q}}
\newcommand{\tfh}{{\widetilde{\mathfrak h}}}
\title[Flagged LLT polynomials, nonsymmetric plethysm, and nonsymme]{Flagged LLT polynomials, nonsymmetric plethysm, and nonsymmetric Macdonald polynomials}
\author{Jonah Blasiak, Mark Haiman, Jennifer Morse, Anna Pun, George H. Seelinger}
\date{Source: arXiv:2506.09015v2 (CC BY 4.0)}
\begin{document}
\maketitle

\noindent\emph{Source and licence.} Flagged LLT polynomials, nonsymmetric plethysm, and nonsymmetric Macdonald polynomials. Version arXiv:2506.09015v2 (CC BY 4.0), distributed under the Creative Commons Attribution 4.0 International licence, as recorded in the arXiv licence metadata for this exact version and shown on the abstract page https://arxiv.org/abs/2506.09015v2. Full rights notice: licenses/arxiv-2506.09015-CC-BY-4.0.txt.\par\smallskip

\noindent\textit{Editorial scope.} Counted unit: the Thm whose source label is \texttt{thm:PiA-symm} in the article (source0.tex lines 2598--2616, source label \texttt{thm:PiA-symm}), delivered together with the article's own complete original proof of it (source0.tex lines 2912--3006, one \texttt{proof} environment of 95 lines). No mathematics is written, repaired or re-derived: statement and proof are the article's own text line for line. The proof is detached from the statement in the source: the article states the result at lines 2598--2616 and prints its proof at lines 2912--3006, and the proof environment's own header names the statement, so the association is the article's own. Required context retained uncounted: e coprod formula (lines 1809--1812); e:htild (lines 1886--1889); prop:symm-by-ha (lines 2860--2888). Every definition, display and label of those lines is retained unchanged. Omitted from this package and not redistributed: 1--1808, 1813--1885, 1890--2597, 2617--2859, 2889--2911, 3007--8812. Nothing inside a retained excerpt is omitted or altered: every retained body is delivered exactly as the article's lines are written, so no omission receipt of any kind accompanies this package. Citations: the retained excerpts carry no citation command at all, so no bibliography entry of the article is transcribed and no reference is redistributed. Reference adaptation: none was needed and none was made; every reference inside the delivered unit resolves inside it, so no unresolved reference or citation is produced. Printed slips kept literal: no printed word, symbol, hypothesis or number of the counted statement or of its proof was altered. Layout and class: the article's own class is \texttt{amsart} with packages including amssymb, amsmath, amsthm, bbm, graphicx, tikz, hyperref, textgreek, stmaryrd, blkarray, geometry; the wrapper is the lane's pinned \texttt{amsart} editorial wrapper at 10pt on A4, which restates the theorem-like environments the retained text needs and adds the article's own macro definitions that the retained text needs (\texttt{\textbackslash NN}, \texttt{\textbackslash aA}, \texttt{\textbackslash bb}, \texttt{\textbackslash fh}, \texttt{\textbackslash pp}, \texttt{\textbackslash qq}, \texttt{\textbackslash tfh}), with extra wrapper packages (bm, bbm, dsfont, xspace, cancel, mathtools). The delivered numbering is the unit's own. Assets: no retained span contains a figure environment, a graphics-inclusion command, a PDF-page inclusion, a table or a TikZ picture; no image file is copied into this package, the package declares no asset dependency and no image file is redistributed with it. Licence: version arXiv:2506.09015v2 of this article is distributed under the Creative Commons Attribution 4.0 International licence, as recorded in the arXiv licence metadata for this exact version and shown on the abstract page https://arxiv.org/abs/2506.09015v2. A full rights notice is delivered at licenses/arxiv-2506.09015-CC-BY-4.0.txt. Characters: every retained excerpt is pure ASCII, so the delivered engine, XeLaTeX with its default Latin Modern text font, reports no missing character.
\begin{equation}
\label{e coprod formula} \fh_{\aA}[X_1+Y_1,\ldots,X_l+Y_l] = \sum_{\pp
+ \qq = \aA} \fh_{\pp}(X_1,\ldots,X_l)\,  \fh_{\qq}(Y_1,\ldots,Y_l) \,,
\end{equation}

\begin{equation}\label{e:htild}
\tfh_{\aA}(X_1,\ldots,X_n) = \fh_{\aA}[X_1,X_2-AX_1, \ldots,
X_n-AX_{n-1}]
\end{equation}

\begin{prop}\label{prop:symm-by-ha}
(a) Given $\aA \in \NN ^{n}$, $1\leq i_{0}\leq i_{1}\leq n$, and
$f(X)\in V_{l}$, suppose that $a_{j} = 0$ for at least $l$ indices
$j\in [i_{0},i_{1}]$.  Then
\begin{equation}\label{e:symm-by-ha}
f[x_{1}+\cdots +x_{i_{1}}] \, \fh_{\aA }[x_{1}, \ldots, x_{n}] = \sum
_{\bb \in \NN ^{i_{1}-i_{0}+1} }c_{\bb }\, \fh_{(\aA';\bb
;\aA''')}[x_{1}, \ldots, x_{n}],
\end{equation}
where $\aA = (\aA '; \aA ''; \aA ''')$ with $\aA' = (a_1,\ldots ,
a_{i_0-1})$, $\aA'' = (a_{i_0},\ldots , a_{i_1})$, and $\aA''' =
(a_{i_1+1}\ldots , a_{n})$, and $c_{\bb }$ depends only on $\aA ''$,
$\bb $, and $f$.

(b) Given $\aA \in \NN ^{n}$, $1\leq i_{0}\leq i_{1}\leq n$, and
$f(X)\in V_{l}$, suppose that $a_{j} = 0$ for at least $l+e$ indices
$j\in [i_{0},i_{1}]$.  Then
\begin{multline}\label{e:symm-by-ha-A}
f[x_{1}+\cdots +x_{i_{1}}] \, \fh_{\aA }[x_{1}, x_{2} - A\, x_{1},
\ldots, x_{n} - A\, x_{n-1}]\\
\equiv \sum _{\bb \in \NN ^{i_{1}-i_{0}+1} } c_{\bb }\, \fh_{(\aA';\bb
;\aA''')}[x_{1}, x_{2} - A\, x_{1}, \ldots, x_{n} - A\, x_{n-1}]
\pmod{\Lambda (A)_{>e}}
\end{multline}
where $\aA = (\aA '; \aA ''; \aA ''')$ with $\aA' = (a_1,\ldots ,
a_{i_0-1})$, $\aA'' = (a_{i_0},\ldots , a_{i_1})$, and $\aA''' =
(a_{i_1+1}\ldots , a_{n})$, and $c_{\bb }$ depends {\rm (mod $\Lambda
(A)_{>e}$)} only on $\aA ''$, $\bb $, and $f$.
\end{prop}

\begin{thm}\label{thm:PiA-symm}
Let $f(X)$ be a symmetric function, and let $g(x) =
g(x_{1},\ldots,x_{r})$ and $h(y) = h(y_{1},\ldots,y_{s})$ be
polynomials, all with coefficients in $\Lambda (A)$.  If $f(X)\in
V_{l}$, $\deg (h)\leq d$, and $n\geq r+l+d+e$, then
\begin{equation}\label{e:PiA-symm-in}
\Pi _{A,(x_{1}, \ldots, x_{n},\, y_{1},\ldots,y_{s})}
\bigl(g(x_{1},\ldots,x_{r})\, f[x_{1}+\cdots +x_{n}] \,
h(y_{1},\ldots,y_{s})\bigr)
\end{equation}
is symmetric {\rm (mod $\Lambda (A)_{>e}$)} in the variables
$x_{r+1},\ldots,x_{n-(l+d+e)+1}$, and upon setting $x_{i} = 0$ for
$n-(l+d+e)+1<i\leq n$, it reduces to
\begin{equation}\label{e:PiA-symm-out}
(\Pi _{A,x}\, g(x))\, f \! \left[\frac{x_{1}+\cdots
+x_{n-(l+d+e)+1}}{1-A} \right]\, (\Pi _{A,y}\, h(y)) \pmod{\Lambda
(A)_{>e}}.
\end{equation}
\end{thm}

\begin{proof}[Proof of Theorem~\ref{thm:PiA-symm}] We set $\tfh _{\aA
}(X_{1},\ldots,X_{k}) = \fh_{\aA }[X_{1}, X_{2} - A\, X_{1}, \ldots,
X_{k} - A\, X_{k-1}]$, as in \eqref{e:htild}.  By linearity, we can
assume that $g(x) = \tfh _{\aA }[x_{1}, \ldots, x_{r}]$ and $h(y) =
\tfh _{\aA '}[y_{1}, \ldots, y_{s}]$, where $\aA \in \NN ^{r}$, $\aA
'\in \NN ^{s}$.  Using the coproduct formula \eqref{e coprod formula},
we can express $\tfh _{\aA '}[y_{1}, \ldots, y_{s}] $ in terms of
functions $\tfh _{\bb }$ of the full list of variables $x, y$, as
\begin{equation}\label{e:hy-A-extended}
\begin{aligned}
\tfh _{\aA '}[y_{1},& \ldots, y_{s}] \\
& = \sum _{\pp +\qq =\aA '} \fh_{\pp }[(A-1)X,0,\ldots,0]\, \fh_{\qq
}[(1-A)X + y_{1}, y_{2} - A\, y_{1}, \ldots, y_{s} - A\, y_{s-1}]\\
& = \sum _{\pp +\qq =\aA '} h_{\pp _{+}}[(A-1)X]\, \tfh
_{(0^{n};\, \qq) }[x_{1}, \ldots, x_{n}, y_{1}, \ldots, y_{s}],
\end{aligned}
\end{equation}
where $X = x_{1}+\cdots +x_{n}$, and $h_{\pp _{+}}$ is the complete
symmetric function indexed by the partition $\pp _{+}$ whose parts are
the nonzero entries of $\pp $.  Then
\begin{equation}\label{e:fgh-by-ha}
g(x)\, f[X]\, h(y) = \sum _{\pp +\qq =\aA '} f[X]\,
h_{\pp _{+}}[(A-1)X]\, \tfh_{(\aA ;\, 0^{n-r};\, \qq) }[x_{1},
\ldots, x_{n}, y_{1}, \ldots, y_{s}].
\end{equation}

Since $V_{d}$ contains all symmetric
functions of degree $\leq d$, we have $f[X]\, h_{\pp _{+}}[(A-1)X] \in
V_{l+d}$.  By Proposition~\ref{prop:symm-by-ha}(b) with
$[i_{0},i_{1}]= [n-(l+d+e)+1,n]$, it follows that
\begin{multline}\label{e:p-term}
f[X]\, h_{\pp _{+}}[(A-1)X]\, \tfh_{(\aA ;\, 0^{n-r};\, \qq)
}[x_{1}, \ldots, x_{n}, y_{1}, \ldots, y_{s}]\\
\equiv \sum _{\bb \in \NN ^{l+d+e}} c_{\pp ,\bb }\, \tfh_{(\aA ;\,
0^{n-r-(l+d+e)};\, \bb ;\, \qq )} [x_{1}, \ldots, x_{n}, y_{1},
\ldots, y_{s}] \pmod{\Lambda (A)_{>e}},
\end{multline}
where the coefficients $c_{\pp ,\bb }\in \Lambda (A)$ are independent
of $n$, $\aA $, and $\qq $.  In particular, by taking $\aA = 0^{r}$
and $\qq = 0^{s}$, we conclude that
\begin{equation}\label{e:cpb-1}
f[X]\, h_{\pp _{+}}[(A-1)X] \equiv \sum _{\bb \in \NN ^{l+d+e}} c_{\pp
,\bb }\, \tfh_{( 0^{n-(l+d+e)};\, \bb )} [x_{1}, \ldots, x_{n}]
\pmod{\Lambda (A)_{>e}},
\end{equation}
for all $n\geq l+d+e$.  Letting $m=n-(l+d+e)$ and setting $x_{m+1} =
\ldots = x_{n} = 0$, this becomes
\begin{equation}\label{e:cpb-2}
f[Z]\, h_{\pp _{+}}[(A-1)Z] \equiv \sum _{\bb \in \NN ^{l+d+e}} c_{\pp
,\bb }\, h _{\bb _{+}} [(1-A)Z] \pmod{\Lambda (A)_{>e}},
\end{equation}
where $Z = x_{1}+\cdots +x_{m}$.  Since \eqref{e:cpb-2} holds for
arbitrarily large $m$, it holds as an identity of symmetric functions
in a formal alphabet $Z$, and hence for plethystic evaluation in any
$Z$.

Now, \eqref{e:fgh-by-ha} and \eqref{e:p-term} imply
\begin{multline}\label{e:PiA-p-term}
\Pi _{A,(x_{1},\ldots,x_{n},y_{1},\ldots,y_{s})}
\bigl( g(x)\, f[x_{1},\ldots,x_{n}]\, h(y) \bigr)\\
\equiv \sum _{\pp +\qq =\aA '}\sum _{\bb \in \NN ^{l+d+e}} c_{\pp ,\bb
}\, \fh_{(\aA ;\, 0^{n-r-(l+d+e)};\, \bb ;\, \qq )} [x_{1}, \ldots,
x_{n}, y_{1}, \ldots, y_{s}] \pmod{\Lambda (A)_{>e}}.
\end{multline}
The block of zeroes in positions $r+1$ through $n-(l+d+e)$ of $(\aA
;\, 0^{n-r-(l+d+e)};\, \bb ;\, \qq )$ implies that each term is
symmetric in $x_{r+1},\ldots,x_{n-(l+d+e)+1}$, as claimed.
Furthermore, on setting $x_{j} = 0$ for $n-(l+d+e)+1 < j \leq n$, the
right hand side of \eqref{e:PiA-p-term} becomes
\begin{equation}\label{e:PiA-p-term-x=0}
\fh_{\aA }[x_{1},\ldots,x_{r}] \sum _{\pp +\qq =\aA '}\sum _{\bb \in \NN
^{l+d+e}} c_{\pp ,\bb }\, h_{\bb_+}[X]\, \fh_{\qq } [X + y_{1}, \ldots,
y_{s}]
\end{equation}
where we now set $X = x_{1}+\cdots +x_{n-(l+d+e)+1}$.  Taking $Z =
X/(1-A)$ in \eqref{e:cpb-2}, we see that \eqref{e:PiA-p-term-x=0} is
congruent (mod $\Lambda (A)_{>e}$) to
\begin{equation}\label{e:PiA-p-term-x=0-bis}
\fh_{\aA }[x_{1},\ldots,x_{r}] f[X/(1-A)] \sum _{\pp +\qq =\aA '} h_{\pp
_{+}}[-X]\, \fh_{\qq } [X + y_{1}, \ldots, y_{s}].
\end{equation}
The summation factor here reduces to
\begin{equation}\label{e:PiA-p-term-x=0-ter}
 \sum _{\pp +\qq =\aA '} \fh_{\pp }[-X,0,\ldots,0]\, \fh _{\qq } [X +
y_{1}, \ldots, y_{s}] = \fh_{\aA '}[y_{1},\ldots,y_{s}],
\end{equation}
so \eqref{e:PiA-p-term-x=0-bis} is equal to
\begin{equation}\label{e:PiA-p-term-x=0-iv}
\fh_{\aA }[x_{1},\ldots,x_{r}] \, f \! \left[\frac{x_{1}+\cdots
+x_{n-(l+d+e)+1}}{1-A} \right] \, \fh_{\aA '}[y_{1},\ldots,y_{s}],
\end{equation}
which is \eqref{e:PiA-symm-out} for $g(x) = \tfh_{\aA
}[x_{1},\ldots,x_{r}]$ and $h(y) = \tfh _{\aA
'}[y_{1},\ldots,y_{s}]$.
\end{proof}

\end{document}
